% Options for packages loaded elsewhere
% Options for packages loaded elsewhere
\PassOptionsToPackage{unicode}{hyperref}
\PassOptionsToPackage{hyphens}{url}
\PassOptionsToPackage{dvipsnames,svgnames,x11names}{xcolor}
\PassOptionsToPackage{space}{xeCJK}
%
%% Load thmtools before ctexbook so \renewtheorem works
\RequirePackage{thmtools}
\documentclass[
  11pt,
  onesidefontset=none,
]{ctexbook}
\usepackage{xcolor}
\usepackage[left=1in,right=1in,top=1in,bottom=1in]{geometry}
\usepackage{amsmath,amssymb}
\setcounter{secnumdepth}{5}
\usepackage{iftex}
\ifPDFTeX
  \usepackage[T1]{fontenc}
  \usepackage[utf8]{inputenc}
  \usepackage{textcomp} % provide euro and other symbols
\else % if luatex or xetex
  \usepackage{unicode-math} % this also loads fontspec
  \defaultfontfeatures{Scale=MatchLowercase}
  \defaultfontfeatures[\rmfamily]{Ligatures=TeX,Scale=1}
\fi
\usepackage{lmodern}
\ifPDFTeX\else
  % xetex/luatex font selection
  \setmainfont[]{TeX Gyre Pagella}
  \ifXeTeX
    \usepackage{xeCJK}
    \setCJKmainfont[]{Kaiti SC}
    \setCJKsansfont[]{Noto Sans SC}
    \setCJKmonofont[]{Noto Sans SC}
  \fi
  \ifLuaTeX
    \usepackage[]{luatexja-fontspec}
    \setmainjfont[]{Kaiti SC}
  \fi
\fi
% Use upquote if available, for straight quotes in verbatim environments
\IfFileExists{upquote.sty}{\usepackage{upquote}}{}
\IfFileExists{microtype.sty}{% use microtype if available
  \usepackage[]{microtype}
  \UseMicrotypeSet[protrusion]{basicmath} % disable protrusion for tt fonts
}{}
\makeatletter
\@ifundefined{KOMAClassName}{% if non-KOMA class
  \IfFileExists{parskip.sty}{%
    \usepackage{parskip}
  }{% else
    \setlength{\parindent}{0pt}
    \setlength{\parskip}{6pt plus 2pt minus 1pt}}
}{% if KOMA class
  \KOMAoptions{parskip=half}}
\makeatother
% Make \paragraph and \subparagraph free-standing
\makeatletter
\ifx\paragraph\undefined\else
  \let\oldparagraph\paragraph
  \renewcommand{\paragraph}{
    \@ifstar
      \xxxParagraphStar
      \xxxParagraphNoStar
  }
  \newcommand{\xxxParagraphStar}[1]{\oldparagraph*{#1}\mbox{}}
  \newcommand{\xxxParagraphNoStar}[1]{\oldparagraph{#1}\mbox{}}
\fi
\ifx\subparagraph\undefined\else
  \let\oldsubparagraph\subparagraph
  \renewcommand{\subparagraph}{
    \@ifstar
      \xxxSubParagraphStar
      \xxxSubParagraphNoStar
  }
  \newcommand{\xxxSubParagraphStar}[1]{\oldsubparagraph*{#1}\mbox{}}
  \newcommand{\xxxSubParagraphNoStar}[1]{\oldsubparagraph{#1}\mbox{}}
\fi
\makeatother


\usepackage{longtable,booktabs,array}
\newcounter{none} % for unnumbered tables
\usepackage{calc} % for calculating minipage widths
% Correct order of tables after \paragraph or \subparagraph
\usepackage{etoolbox}
\makeatletter
\patchcmd\longtable{\par}{\if@noskipsec\mbox{}\fi\par}{}{}
\makeatother
% Allow footnotes in longtable head/foot
\IfFileExists{footnotehyper.sty}{\usepackage{footnotehyper}}{\usepackage{footnote}}
\makesavenoteenv{longtable}
\usepackage{graphicx}
\makeatletter
\newsavebox\pandoc@box
\newcommand*\pandocbounded[1]{% scales image to fit in text height/width
  \sbox\pandoc@box{#1}%
  \Gscale@div\@tempa{\textheight}{\dimexpr\ht\pandoc@box+\dp\pandoc@box\relax}%
  \Gscale@div\@tempb{\linewidth}{\wd\pandoc@box}%
  \ifdim\@tempb\p@<\@tempa\p@\let\@tempa\@tempb\fi% select the smaller of both
  \ifdim\@tempa\p@<\p@\scalebox{\@tempa}{\usebox\pandoc@box}%
  \else\usebox{\pandoc@box}%
  \fi%
}
% Set default figure placement to htbp
\def\fps@figure{htbp}
\makeatother



\ifLuaTeX
\usepackage[bidi=basic,shorthands=off,provide=*]{babel}
\else
\usepackage[bidi=default,shorthands=off,provide=*]{babel}
\fi
\ifPDFTeX
\else
\babelfont{rm}[]{TeX Gyre Pagella}
\fi
\ifLuaTeX
  \usepackage{selnolig} % disable illegal ligatures
\fi


\setlength{\emergencystretch}{3em} % prevent overfull lines

\providecommand{\tightlist}{%
  \setlength{\itemsep}{0pt}\setlength{\parskip}{0pt}}



 
\usepackage[]{natbib}
\bibliographystyle{plainnat-doi}


%% ============================================================
%% Preamble: font, layout, and style settings
%% Matching the style of article.tex (XeLaTeX, Palatino-like)
%% ============================================================

%% ---- Load thmtools early (before amsthm) to enable \renewtheorem ----
\RequirePackage{thmtools}
\usepackage{xpatch}

%% ---- natbib citation style (numbers, sorted & compressed) ----
\setcitestyle{numbers,sort&compress,square}

%% ---- Line spacing ----
\renewcommand{\baselinestretch}{1.2}
%%
\usepackage{thmtools}
\usepackage{orcidlink}
%% ---- Caption: disable auto-centering of single-line captions ----
\usepackage{caption}
\captionsetup{singlelinecheck=off}

%% ---- Custom differential operator ----
\newcommand*{\dd}{\mathop{}\!\mathrm{d}}

%% ---- Allow display math to break across pages ----
\allowdisplaybreaks

%% ---- Hyperlink colors (deferred until hyperref is loaded) ----
\AtBeginDocument{%
  \renewcommand{\contentsname}{目录}%
  \hypersetup{%
    bookmarksnumbered=true,
    bookmarksopen=true,
    colorlinks=true,
    linkcolor=red,
    anchorcolor=blue,
    citecolor=blue,
    breaklinks=true
  }%
}

%% ---- Page style: header (section title centered) + footer (page number centered) ----
\usepackage{fancyhdr}
\usepackage{emptypage}
\pagestyle{fancy}
\fancyhf{}
\fancyhead[C]{\rightmark}
\fancyfoot[C]{\thepage}
\renewcommand{\headrulewidth}{0pt}
\renewcommand{\footrulewidth}{0pt}
\fancypagestyle{plain}{%
  \fancyhf{}%
  \fancyhead[C]{\rightmark}%
  \fancyfoot[C]{\thepage}%
  \renewcommand{\headrulewidth}{0pt}%
  \renewcommand{\footrulewidth}{0pt}%
}
\renewcommand{\sectionmark}[1]{\markright{\thesection\space #1}}

%% ---- Font setup (XeLaTeX, Palatino-like: TeX Gyre Pagella) ----
%% Text font: newpxtext gives Palatino-like (TeX Gyre Pagella) text
%% Math font: Quarto's default unicode-math (Latin Modern Math)
\defaultfontfeatures{Mapping=tex-text}
\XeTeXlinebreaklocale "zh"
\XeTeXlinebreakskip = 0pt plus 1pt minus 0.1pt
\usepackage[theoremfont]{newpxtext}

%% ---- CJK font setup (ctexbook with fontset=none) ----
%%\setCJKmainfont{KaiTi}
%%\setCJKsansfont{Noto Sans SC}
%%\setCJKmonofont{Noto Sans SC}
\punctstyle{banjiao}

%------------------------------------------------------------------------------
% 17. Equation reference parentheses — Quarto 对公式交叉引用使用 \ref
%     而非 \eqref，导致编号不带括号。AtBeginDocument 在 hyperref 加载后执行，
%     确保 \ref 重定义生效（覆盖 hyperref 版本）。
%------------------------------------------------------------------------------
\usepackage{xstring}
\AtBeginDocument{%
  \let\quartoOldRef\ref
  \renewcommand{\ref}[1]{%
    \IfBeginWith{#1}{eq-}{%
      \textup{(\quartoOldRef{#1})}%
    }{%
      \quartoOldRef{#1}%
    }%
  }%
}
% Theorem boxes — rounded corners, white/transparent background
\usepackage[most,theorems,skins,breakable]{tcolorbox}

% --- Color definitions ---
\definecolor{thmframe}{RGB}{120,185,225}

% --- Shared tcolorbox style for theorem-like environments ---
\tcbset{
  thmboxstyle/.style={
    enhanced jigsaw,
    colback=white,
    colframe=thmframe,
    boxrule=0.6pt,
    arc=4pt,
    before skip=3pt plus 4pt minus 2pt,
    after skip=3pt plus 4pt minus 2pt,
    breakable,
    left=1pt,
    right=1pt,
    top=1pt,
    bottom=1pt,
  }
}

\tcolorboxenvironment{theorem}{thmboxstyle}{}{}
\tcolorboxenvironment{lemma}{thmboxstyle}{}{}
\tcolorboxenvironment{corollary}{thmboxstyle}{}{}
\tcolorboxenvironment{proposition}{thmboxstyle}{}{}
\tcolorboxenvironment{definition}{thmboxstyle}{}{}
\tcolorboxenvironment{example}{thmboxstyle}{}{}
\tcolorboxenvironment{remark}{thmboxstyle}{}{}
\makeatletter
\@ifpackageloaded{caption}{}{\usepackage{caption}}
\AtBeginDocument{%
\ifdefined\contentsname
  \renewcommand*\contentsname{目录}
\else
  \newcommand\contentsname{目录}
\fi
\ifdefined\listfigurename
  \renewcommand*\listfigurename{图索引}
\else
  \newcommand\listfigurename{图索引}
\fi
\ifdefined\listtablename
  \renewcommand*\listtablename{表索引}
\else
  \newcommand\listtablename{表索引}
\fi
\ifdefined\figurename
  \renewcommand*\figurename{图}
\else
  \newcommand\figurename{图}
\fi
\ifdefined\tablename
  \renewcommand*\tablename{表}
\else
  \newcommand\tablename{表}
\fi
}
\@ifpackageloaded{float}{}{\usepackage{float}}
\floatstyle{ruled}
\@ifundefined{c@chapter}{\newfloat{codelisting}{h}{lop}}{\newfloat{codelisting}{h}{lop}[chapter]}
\floatname{codelisting}{列表}
\newcommand*\listoflistings{\listof{codelisting}{列表索引}}
\usepackage{amsthm}
\theoremstyle{plain}
\newtheorem{theorem}{定理}[section]
\theoremstyle{plain}
\newtheorem{corollary}{推论}[section]
\theoremstyle{definition}
\newtheorem{definition}{定义}[section]
\theoremstyle{plain}
\newtheorem{lemma}{引理}[section]
\theoremstyle{remark}
\AtBeginDocument{\renewcommand*{\proofname}{证明}}
\newtheorem*{remark}{注记}
\newtheorem*{solution}{解}
\newtheorem{refremark}{注记}[section]
\newtheorem{refsolution}{解}[section]
\makeatother
\makeatletter
\makeatother
\makeatletter
\@ifpackageloaded{caption}{}{\usepackage{caption}}
\@ifpackageloaded{subcaption}{}{\usepackage{subcaption}}
\makeatother
\usepackage{bookmark}
\IfFileExists{xurl.sty}{\usepackage{xurl}}{} % add URL line breaks if available
\urlstyle{same}
\hypersetup{
  pdftitle={基于 Quarto 与标准 LaTeX Book 文档类的中文书籍模板},
  pdfauthor={你的名字; John Doe},
  pdflang={zh-CN},
  pdfkeywords={Quarto, LaTeX, book, Pandoc, template},
  colorlinks=true,
  linkcolor={blue},
  filecolor={Maroon},
  citecolor={Blue},
  urlcolor={Blue},
  pdfcreator={LaTeX via pandoc}}


\title{\Large 基于 Quarto 与标准 LaTeX Book 文档类的中文书籍模板}

\author{\href{https://your-university.edu/your-profile}{你的名字}\footnote{通讯作者: 你的大学 你的院系, 你的城市, 0, 中国. E-mail: \texttt{your.email@university.edu}}\,~\orcidlink{0000-0000-0000-0001}\hspace{2em}John
Doe\footnote{Anonymous University}}

\date{\small 2023-06-23}

\begin{document}
\maketitle
\clearpage
\mbox{}\thispagestyle{empty}
\clearpage
\pagenumbering{Roman}
\markright{}

% --- Chinese theorem/proof names (redefine via \let + \newtheorem) ---
\makeatletter
\let\theorem\relax\let\c@theorem\relax
\let\lemma\relax\let\c@lemma\relax
\let\corollary\relax\let\c@corollary\relax
\let\definition\relax\let\c@definition\relax
\let\proposition\relax\let\c@proposition\relax
\let\example\relax\let\c@example\relax
\let\remark\relax\let\c@remark\relax
\makeatother
\newtheorem{theorem}{定理}[section]
\newtheorem{lemma}{引理}[section]
\newtheorem{corollary}{推论}[section]
\newtheorem{definition}{定义}[section]
\newtheorem{proposition}{命题}[section]
\newtheorem{example}{例}[section]
\newtheorem{remark}{注}[section]
\renewcommand{\proofname}{证明}
\renewcommand*\contentsname{目录}
{
\setcounter{tocdepth}{1}
\tableofcontents
}

\mainmatter
\chapter{Introduction}\label{sec-intro}

This file is a sample book demonstrating Quarto's capabilities with the
standard LaTeX \texttt{book} document class. It supports PDF (via LaTeX)
output. It is assumed that the reader has a basic working knowledge of
stochastic calculus \citet{karatzas1991brownian} and stochastic control
theory \citet{yong1999sch}.

\[
\begin{aligned}
&\forall x \in \mathbb{R},\quad 
\exists\, y \in \mathbb{N}:\quad 
\alpha + \beta = \gamma \neq \delta \leq \epsilon \approx \zeta \\
&\sum_{i=1}^{n} x_i^2 + \prod_{j=1}^{m} y_j^2 = \int_{a}^{b} f(x)\,dx \\
&\lim_{n\to\infty} \frac{1}{n} = 0, \qquad
\nabla \cdot \vec{E} = \frac{\rho}{\varepsilon_0} \\
&A = \begin{pmatrix} a_{11} & a_{12} \\ a_{21} & a_{22} \end{pmatrix}, \qquad
B = \begin{bmatrix} b_1 \\ b_2 \end{bmatrix} \\
&S = \{ x \mid x > 0 \}, \qquad
f: X \to Y, \qquad
x \mapsto f(x) \\
&\bigcup_{i=1}^n A_i \cap \bigcap_{j=1}^m B_j, \qquad
A \subseteq B, \quad A \supset B \\
&\mathcal{L} = \mathcal{F}(\mathbb{R}), \qquad
\mathfrak{g} \cong \mathfrak{h} \\
&\boxed{E = mc^2}
\end{aligned}
\]

\begin{theorem}[Residue
Theorem]\protect\hypertarget{thm-residue}{}\label{thm-residue}

Let \(f\) be analytic in a simply connected domain \(D\). Then for any
closed contour \(\gamma \subset D\),
\begin{equation}\protect\phantomsection\label{eq-einstein}{
\oint_\gamma f(z)\,dz = 2\pi i \sum_{k} \mathrm{Res}(f, z_k)
}\end{equation}

\end{theorem}

This template supports standard Quarto cross-references. You can refer
to Theorem 定理~\ref{thm-residue} in the text.

\chapter{Markdown Basics}\label{sec-markdown}

The template is built on Quarto and Pandoc, \ref{eq-einstein} leveraging
standard LaTeX tools for PDF output.

\section{Equations}\label{equations}

Inline math: \(E = mc^2\). Display math:

\[
\frac{\partial u}{\partial t} = \alpha \nabla^2 u
\]

\section{Theorems and Proofs}\label{theorems-and-proofs}

\begin{lemma}[Key Lemma]\protect\hypertarget{lem-key}{}\label{lem-key}

If \(f\) is continuous on \([a,b]\), then \(f\) is integrable on
\([a,b]\).

\end{lemma}

\begin{corollary}[Easy
Corollary]\protect\hypertarget{cor-easy}{}\label{cor-easy}

Every continuous function on a closed interval is bounded.

\end{corollary}

\begin{definition}[Model
Definition]\protect\hypertarget{def-model}{}\label{def-model}

A stochastic process \(\{X_t\}_{t \ge 0}\) is a collection of random
variables indexed by time.

\end{definition}

\section{Lists}\label{lists}

\begin{enumerate}
\def\labelenumi{\arabic{enumi}.}
\tightlist
\item
  First item
\item
  Second item
\item
  Third item
\end{enumerate}

\begin{itemize}
\tightlist
\item
  Bullet one
\item
  Bullet two
\end{itemize}

\section{Tables}\label{tables}

\begin{longtable}[]{@{}lll@{}}
\caption{A sample table}\label{tbl-sample}\tabularnewline
\toprule\noalign{}
Variable & Description & Unit \\
\midrule\noalign{}
\endfirsthead
\toprule\noalign{}
Variable & Description & Unit \\
\midrule\noalign{}
\endhead
\bottomrule\noalign{}
\endlastfoot
\(x\) & Position & m \\
\(v\) & Velocity & m/s \\
\(a\) & Acceleration & m/s² \\
\end{longtable}

\chapter{参考文献设置}\label{ux53c2ux8003ux6587ux732eux8bbeux7f6e}

This template uses \texttt{natbib} with the \texttt{plainnat-doi}
bibliography style for author-year citations
\citet{karatzas1991brownian}.

\chapter*{致谢}\label{ux81f4ux8c22}
\addcontentsline{toc}{chapter}{致谢}

This should be a simple paragraph before the References to thank those
individuals and institutions who have supported your work on this book.

{}

Use the \texttt{{[}{]}\{.appendix\ options="An\ Appendix"\}} markup if
you have a single appendix. Do not use
\texttt{\textbackslash{}section\{\}} after
\texttt{\textbackslash{}appendix} in LaTeX.

% Bibliography will be placed here by Pandoc
% Patch \bibliography to use \bibname (Chinese "参考文献")
\xpatchcmd{\bibliography}{\chapter{Bibliography}}{\chapter{\bibname}}{}{}
\let\stdchaptermark\chaptermark
\renewcommand{\chaptermark}[1]{}
\makeatletter
\let\stdmkboth\@mkboth
\def\@mkboth#1#2{}
\makeatother
\markright{}
\backmatter
\renewcommand\bibname{参考文献}
\bibliography{bibliography.bib}


\vspace*{18pt}
\noindent\textbf{你的名字}\par

Author biography text goes here.\par
\vspace*{18pt}
\noindent\textbf{John Doe}\par

John Doe is a researcher at Anonymous University.\par

\end{document}
